\section{Design}
\label{sec:design}

\subsection{Core grammar}
\label{subsec:grammar}
\todo{Lead with a single concrete example (s415, cheating) shown in full
Yuho syntax, then unpack each block.}

\begin{quote}
\begin{small}
\texttt{statute 415 "Cheating" \{ \\
~~jurisdiction := SG; \\
~~effective 1872-01-01; \\
~~elements \{ \\
~~~~deception : actus\_reus := "by deceiving any person"; \\
~~~~induction  : actus\_reus := "fraudulently induces" causedBy deception; \\
~~~~mens       : mens\_rea  := "intent to defraud"; \\
~~\} \\
~~penalty alternative \{ \\
~~~~imprisonment := up\_to 3 years; \\
~~~~fine := unlimited; \\
~~~~or\_both; \\
~~\} \\
~~illustration "(a) ..."; \\
~~exception "consent" guard "victim consents" priority 1; \\
\}}
\end{small}
\end{quote}

\subsection{Element decomposition}
\label{subsec:elements}
Every offence is a conjunction (or carefully-named disjunction) of typed
elements. The element kinds --- \texttt{actus\_reus}, \texttt{mens\_rea},
\texttt{circumstance}, \texttt{obligation}, \texttt{prohibition},
\texttt{permission} --- are the type system. Burden qualifiers
(\texttt{burden prosecution}, \texttt{burden defence}) and proof standards
attach per-element. Causal links (\texttt{causedBy other\_element}) and
temporal relations (\texttt{precedes}, \texttt{during}, \texttt{after})
make element graphs traversable.

\subsection{Penalty combinators}
\label{subsec:penalty}
\todo{Diagram: penalty AST tree for s302 (``death or imprisonment for life,
and shall also be liable to fine'') showing nested combinator from G12.}
Three top-level combinators (\texttt{cumulative}, \texttt{alternative},
\texttt{or\_both}); a conditional branch combinator (\texttt{when <ident>});
and a nested-combinator escape hatch added in gap \textsf{G12} for
penalties of the form ``A and (B or C)''. Sentinels \texttt{fine := unlimited}
(\textsf{G8}) and \texttt{caning := unspecified} (\textsf{G14}) cover
statute-text-says-no-cap cases without forcing fabricated numeric ranges.

\subsection{Exceptions as Catala-style defaults}
\label{subsec:exceptions}
\todo{Show the priority DAG; mention how it compiles to Z3 prioritised
rewriting.}
General exceptions (Chapter IV) and section-local provisos are encoded
as priority-ordered \texttt{exception} blocks with an optional
\texttt{defeats} edge to the elements they override. Compilation to Z3
follows Catala's prioritised-rewriting semantics: lower-priority rules
fire only when no higher-priority rule applies.

\subsection{Subsections and amendment lineage}
A statute can carry nested \texttt{subsection (N) \{ ... \}} blocks
(gap \textsf{G5}) with their own elements, penalties, and exceptions.
Multiple \texttt{effective <date>} clauses (gap \textsf{G6}) record the
amendment history; a \texttt{repealed <date>} clause closes the section.

\subsection{The fourteen grammar gaps}
\label{subsec:gaps}
\todo{Compact table summarising G1--G14, status (fixed / deferred / not-a-gap),
and the artefact that resolved each. Source: \texttt{todo.md} Phase D table.}

Mass-encoding the Penal Code surfaced fourteen distinct grammar shortcomings.
We classify each as a real grammar limitation (10 items) or a not-a-gap
that the diagnostic layer should catch (4 items). All ten real gaps
were resolved in the grammar; the four diagnostic-layer items became
fidelity linters (\S\ref{subsec:fidelity}).
